\section{Dimension and cohomology of the boundary}
\label{sec:boundary-topology}

The compactification $X=\Gamma\sqcup B$ has two topological properties
that will drive the contradiction: $B$ has covering dimension at most one,
and its first \v{C}ech cohomology vanishes.  Neither assertion requires
local connectivity.  The first comes from the small accumulation sets of
coordinate cuts; the second transfers a finite-cover cocycle to the
presentation complex and removes its finite coboundary there.

\subsection{Thin transition sets and covering dimension}

For a coordinate $F$ and a level $t$ avoiding its vertex values, let
$P(F,t)$ be the set of endpoints of edges in $S(F,t)$, and put
\[
 E(F,t)=\overline{P(F,t)}\cap B.
\]
Thus $E(F,t)$ records where the cut accumulates in the boundary.
All closures and diameters below use the metric $\rho$ on $X$.

\begin{lemma}\label{topo:thin-walls}
For every coordinate $F$, the compact set $E(F,t)$ has
$\mathcal H^1_\rho$-measure zero for almost every $t\in(0,1)$.
\end{lemma}

\begin{proof}
First suppose that $F$ is a base coordinate.  The level-length estimate
of Lemma~\ref{dyn:length}, with the identity translate, gives
\[
 \sum_{e\in S(F,t)}\ \sum_{e'\text{ near }e}\lambda(e')<\infty
\]
for almost every $t$.  Here the fixed neighborhood of an edge includes
the boundaries of its incident cells.  Traversing these boundaries along
a simple dual path shows that the diameter of the endpoint set of a dual
component $D$ is at most
\[
 C\sum_{e\in D}\ \sum_{e'\text{ near }e}\lambda(e'),
\]
where $C$ depends only on the finite presentation.  Bounded local
incidence bounds the multiplicity with which a nearby edge is used.

Fix a level with the finite sum above.  The wall has finitely many dual
components, and the dual adjacency graph is locally finite.  Delete
finite subsets $D_n\subset S(F,t)$ increasing to the entire wall.  At
each stage only finitely many dual components remain: deleting finitely
many vertices from a connected locally finite graph creates only finitely
many components.  Their endpoint-set closures cover $E(F,t)$, since
deleting finitely many isolated graph vertices does not change boundary
accumulation.  The sum of their diameters is at most
\[
 C\sum_{e\in S(F,t)\setminus D_n}
       \ \sum_{e'\text{ near }e}\lambda(e')\longrightarrow0.
\]
Their maximum diameter also tends to zero.  These finite covers prove
that $\mathcal H^1_\rho(E(F,t))=0$.

For a localized coordinate, its straddling edges are contained in a
finite union of base-coordinate walls at the shifted levels appearing
in its definition.  The corresponding boundary accumulation set is
therefore contained in the finite union of their accumulation sets.
Each shifted level is admissible for almost every original level;
translations and reflections of the level parameter preserve null sets.
The result follows.
\end{proof}

\begin{proposition}\label{topo:dimension}
The boundary $B$ has covering dimension at most one.
\end{proposition}

\begin{proof}
We construct arbitrarily fine finite open covers of multiplicity at most
two.  Choose a finite grid of levels in finitely many coordinates.  Each
level avoids vertex values and satisfies Lemma~\ref{topo:thin-walls}.
Let $E$ be the union of the corresponding sets $E(F,t)$.  It is compact
and has $\mathcal H^1_\rho$-measure zero.

The set $E$ has arbitrarily small finite clopen partitions.  Indeed, for
$x\in E$, the distance map $y\mapsto\rho(x,y)$ is Lipschitz, so its image
on $E$ has Lebesgue measure zero.  There are arbitrarily small positive
radii $r$ for which the sphere of radius $r$ about $x$ misses $E$.
The sets $E\cap B_\rho(x,r)$ are then clopen in $E$.  Compactness gives a
finite small clopen cover, which becomes a disjoint partition by
successively subtracting its previous members.

Each selected cut has a locally constant sign on the complement of its
accumulation set.  To define it at $x\notin E(F,t)$, choose a neighborhood
$W$ of $x$ in $X$ that contains no endpoint of a straddling edge.
Lemma~\ref{compact:near-path} joins all vertices sufficiently near $x$
by paths with vertices in $W$.  Along such a path the sign of $F-t$
cannot change.  Hence all sufficiently near vertices have the same sign.
This sign also applies at all sufficiently nearby boundary points,
so it is locally constant.  This definition is valid even when
$F(x)=t$: it uses the signs on graph vertices approaching $x$.

The finitely many sign tuples partition $B\setminus E$ into finitely many
disjoint sets open in $B$.  Within any one such set, each tested
coordinate belongs to the closure of one grid interval.  First choose
enough coordinates that the remaining tail of the metric
$\rho=\sum_j c_j|F_j(x)-F_j(y)|$ is small, and then choose the grids sufficiently
fine.  Every sign-tuple set consequently has arbitrarily small diameter.

Partition $E$ into finitely many small clopen sets.  They are disjoint
compact subsets of $B$, so they admit pairwise disjoint open neighborhoods
in $B$, still of arbitrarily small diameter.  These neighborhoods,
together with the sign-tuple sets, cover $B$.  At each point at most one
neighborhood and at most one sign-tuple set occur.  The cover therefore
has multiplicity at most two.  A Lebesgue number for any prescribed
finite open cover shows that a sufficiently fine cover of this form
refines it.  This is the assertion $\dim B\le1$.
\end{proof}

\subsection{A finite-support correction and \v{C}ech cohomology}

The next argument uses only the near-path property, the shrinking
diameters of cell boundaries at infinity, and the group-cohomological
vanishing from Lemma~\ref{found:duality}.  In particular, it does
not replace \v{C}ech cohomology by singular cohomology on the possibly
nonlocally connected boundary.

\begin{proposition}\label{topo:cech}
With coefficients $k=\mathbb F_2$, one has $\check H^1(B;k)=0$.
\end{proposition}

\begin{proof}
Let $d=(d_{ij})$ be a simplicial $1$-cocycle on the nerve of a finite open
cover $\{U_i\}_{i=1}^N$ of $B$, with $d_{ii}=0$.  We show that it becomes a
coboundary after refinement.

\paragraph{Labeling graph vertices.}
Choose $\delta>0$ so that every ball of radius $3\delta$ in $B$ is
contained in some $U_i$.  For each vertex $u\in\Gamma$, choose a nearest
point $b(u)\in B$ and a label $i(u)$ with
\[
 B_B(b(u),3\delta)\subset U_{i(u)}.
\]
If $u\to x\in B$, then
$\rho(b(u),x)\le2\rho(u,x)\to0$.  Thus all labels occurring sufficiently
near $x$ contain a fixed neighborhood of $x$ in $B$.

Define a cellular $1$-cochain on $\mathcal C$ by
\[
 \alpha(uv)=d_{i(u)i(v)}
\]
when that pair occurs in the cover nerve, and by zero otherwise.
Its coboundary has finite support.  Otherwise choose infinitely many
distinct cells where the coboundary is nonzero.  Local finiteness forces
their vertices to escape every finite set.  Each cell boundary has at
most three edges, whose lengths tend to zero, so a subsequence of the
vertex sets converges to one point $x\in B$.  All labels on every
sufficiently late cell contain $x$.  The nerve cocycle identity then
gives zero around its boundary, a contradiction.  Repeated labels or
boundary edges cause no difficulty: the identity is a cellular sum,
and $d_{ii}=0$.

\paragraph{The group-ring cocycle.}
We explain why the finite-support cochain $\omega=d\alpha$ represents
a group $2$-cocycle.  Start a free $k\Gamma$-resolution with the cellular
chains of $\mathcal C$ through degree two.  This is possible because
$\mathcal C$ is simply connected: choose the next free module to map
onto $\ker\partial_2$, and continue.  For any finite-support scalar
cellular cochain $\omega$, the formula
\[
 \Phi_\omega(c)=\sum_{g\in\Gamma}\omega(g^{-1}c)g
\]
defines a $k\Gamma$-module cochain.  The sum is finite for every finite
cellular chain $c$, and its identity coefficient is $\omega(c)$.
Conversely, since the quotient has finitely many cells in each relevant
degree, every module cochain through degree two is obtained in this way.

For every finite cellular $2$-cycle $z$ and every $g\in\Gamma$,
\[
 \omega(g^{-1}z)=\alpha\bigl(\partial(g^{-1}z)\bigr)=0.
\]
Thus $\Phi_\omega$ vanishes on $\ker\partial_2$, so it is a cocycle on
the extended free resolution.  Lemma~\ref{found:duality} gives
$H^2(\Gamma,k\Gamma)=0$.  There is therefore a module $1$-cochain whose
coboundary is $\Phi_\omega$.  Taking identity coefficients yields a
finite-support scalar $1$-cochain $\eta$ with
\[
 d\eta=d\alpha.
\]
Ordinary simple connectivity of $\mathcal C$ now supplies a scalar
function $l$ on its vertices such that
\[
 \alpha-\eta=dl.
\]

\paragraph{Returning to the boundary.}
Fix $x\in U_i$.  Choose a neighborhood $W$ of $x$ in $X$ avoiding all
endpoints of edges in the support of $\eta$, and small enough that every
vertex label occurring in $W$ contains $x$.  If $uv$ is an edge with
vertices in $W$, the sets $U_i,U_{i(u)},U_{i(v)}$ meet at $x$.
The cocycle identity and the equation $dl=\alpha$ on this edge imply
\[
 l(u)+d_{i,i(u)}=l(v)+d_{i,i(v)}.
\]
By Lemma~\ref{compact:near-path}, this expression has one constant value
on all vertices sufficiently near $x$.  Denote that value by $l_i(x)$.
It is independent of the chosen neighborhood, since graph vertices
approach every boundary point.

The function $l_i$ is locally constant on $U_i$.  Indeed, a smaller
neighborhood on whose graph vertices the displayed expression is
constant also computes the same germ at every sufficiently nearby
boundary point.  On an overlap $U_i\cap U_j$, the cocycle identity gives
\[
 l_i+l_j=d_{ij}.
\]
Refine the cover by the open sets
$U_i\cap l_i^{-1}(a)$, for $a\in k$, omitting empty sets.  On this finite
refinement the original cocycle is the coboundary of the assignment $a$
to the corresponding cover member.  Its \v{C}ech class vanishes.
\end{proof}

We have proved that $B$ has covering dimension at most one and
has no first \v{C}ech cohomology.  These properties alone do not ensure
local connectivity.  The next steps isolate a continuum without cut
points and prove its local connectivity by pushing graph paths toward
the boundary.
